\documentclass[10pt,oneside]{article}
\usepackage[a4paper,margin=20mm,headheight=15pt]{geometry}
\usepackage[T1]{fontenc}
\usepackage[utf8]{inputenc}
\usepackage{lmodern}
\usepackage{amsmath,amssymb,mathtools}
\usepackage{xcolor}
\usepackage{listings}
\usepackage{enumitem}
\usepackage{fancyhdr}
\usepackage{hyperref}
\usepackage{microtype}
\definecolor{codebg}{HTML}{F6F7F9}
\definecolor{codeframe}{HTML}{CBD2D9}
\definecolor{accent}{HTML}{284B63}
\hypersetup{colorlinks=true,linkcolor=accent,urlcolor=accent}
\lstdefinelanguage{Isabelle}{morekeywords={theory,imports,begin,end,type_synonym,definition,where,theorem,shows,locale,for,fixes,assumes},sensitive=true,morecomment=[s]{(*}{*)},morestring=[b]"}
\lstset{language=Isabelle,basicstyle=\ttfamily\scriptsize,backgroundcolor=\color{codebg},frame=single,rulecolor=\color{codeframe},breaklines=true,breakatwhitespace=false,columns=fullflexible,keepspaces=true,showstringspaces=false,xleftmargin=0.6em,xrightmargin=0.3em}
\setlist{nosep,leftmargin=2em}
\setcounter{tocdepth}{2}
\setcounter{secnumdepth}{2}
\pagestyle{fancy}
\fancyhf{}
\fancyhead[L]{Translation-review dossier}
\fancyhead[R]{Translation B}
\fancyfoot[C]{\thepage}
\newenvironment{sourcequote}{\begin{quote}\itshape}{\end{quote}}
\newcommand{\reviewlabel}[1]{\par\smallskip\noindent\textsf{\bfseries #1}\par\smallskip}
\emergencystretch=3em
\sloppy
\begin{document}
\begin{titlepage}
\centering
\vspace*{0.15\textheight}
{\LARGE\bfseries Translation-review dossier\par}
\vspace{1.2em}
{\Large Conditional Alessandrini uniqueness for planar \(L^p\) potentials\par}
\vspace{1em}
{\large Independent translation B\par}
\vfill
{Package version 2\par}
{Human-ratification state: pending\par}
\end{titlepage}
\tableofcontents
\clearpage

\section{Scope and trust status}
This dossier covers the frozen conditional v002 target, not the manuscript's
Dirichlet-to-Neumann theorem. The proof is relative to the ten literature
interfaces reproduced below. The active approval mode is
\texttt{testing\_waiver}; exact-byte human ratification of the statement and
all literature entries remains pending.

\section{Definitions}
The following byte-pinned project definitions are the complete meaning-bearing
closure of the advertised v002 proposition, in dependency order.
% CF-DEFINITIONS-BEGIN
% CF-DEFINITION-ENTRY-BEGIN: DEF-01-TYPES
\subsubsection{\texttt{\detokenize{DEF-01-TYPES}}}
\noindent\textit{Formal identifier:} \texttt{\detokenize{slp_point; slp_scalar_field; slp_gradient_field; slp_h1_data}}

\reviewlabel{Isabelle code}
% CF-ISABELLE-BEGIN: DEF-01-TYPES
\begin{lstlisting}
type_synonym slp_point = "real ^ 2"
type_synonym slp_scalar_field = "slp_point \<Rightarrow> complex"
type_synonym slp_gradient_field = "slp_point \<Rightarrow> complex ^ 2"
type_synonym slp_h1_data = "slp_scalar_field \<times> slp_gradient_field"
\end{lstlisting}
% CF-ISABELLE-END: DEF-01-TYPES

\reviewlabel{Independent English translation of the Isabelle code}
% CF-BACK-TRANSLATION-BEGIN: DEF-01-TYPES
The formal universe identifies the plane with $\mathbb{R}^2$. Complex functions on it are the scalar representatives; their recorded weak gradients are complex $2$-vectors, and an $H^1$ datum stores both representatives as an ordered pair.
% CF-BACK-TRANSLATION-END: DEF-01-TYPES
% CF-DEFINITION-ENTRY-END: DEF-01-TYPES

% CF-DEFINITION-ENTRY-BEGIN: DEF-02-SMOOTH-BOUNDARY
\subsubsection{\texttt{\detokenize{DEF-02-SMOOTH-BOUNDARY}}}
\noindent\textit{Formal identifier:} \texttt{\detokenize{slp_smooth_boundary_at}}

\reviewlabel{Isabelle code}
% CF-ISABELLE-BEGIN: DEF-02-SMOOTH-BOUNDARY
\begin{lstlisting}
definition slp_smooth_boundary_at ::
  "slp_point set \<Rightarrow> slp_point \<Rightarrow> bool"
where
  "slp_smooth_boundary_at D x \<longleftrightarrow>
    (\<exists>(U :: slp_point set) (W :: slp_point set)
       (f :: slp_point \<Rightarrow> slp_point)
       (g :: slp_point \<Rightarrow> slp_point) (i :: 2).
      open U \<and> open W \<and> x \<in> U \<and>
      homeomorphism U W f g \<and>
      smooth_on U f \<and> smooth_on W g \<and>
      f ` (D \<inter> U) = W \<inter> {y. 0 < y $ i})"
\end{lstlisting}
% CF-ISABELLE-END: DEF-02-SMOOTH-BOUNDARY

\reviewlabel{Independent English translation of the Isabelle code}
% CF-BACK-TRANSLATION-BEGIN: DEF-02-SMOOTH-BOUNDARY
At $x$, the set $D$ must admit a smooth local flattening: on open neighborhoods, a smooth homeomorphism with smooth inverse sends $D$ exactly to the part where one selected coordinate is positive.
% CF-BACK-TRANSLATION-END: DEF-02-SMOOTH-BOUNDARY
% CF-DEFINITION-ENTRY-END: DEF-02-SMOOTH-BOUNDARY

% CF-DEFINITION-ENTRY-BEGIN: DEF-03-DOMAIN
\subsubsection{\texttt{\detokenize{DEF-03-DOMAIN}}}
\noindent\textit{Formal identifier:} \texttt{\detokenize{slp_bounded_smooth_domain}}

\reviewlabel{Isabelle code}
% CF-ISABELLE-BEGIN: DEF-03-DOMAIN
\begin{lstlisting}
definition slp_bounded_smooth_domain :: "slp_point set \<Rightarrow> bool"
where
  "slp_bounded_smooth_domain D \<longleftrightarrow>
    D \<noteq> {} \<and> open D \<and> connected D \<and> bounded D \<and>
    (\<forall>x\<in>frontier D. slp_smooth_boundary_at D x)"
\end{lstlisting}
% CF-ISABELLE-END: DEF-03-DOMAIN

\reviewlabel{Independent English translation of the Isabelle code}
% CF-BACK-TRANSLATION-BEGIN: DEF-03-DOMAIN
The predicate selects nonempty bounded connected open planar sets whose boundary is locally smoothly flattenable at each point.
% CF-BACK-TRANSLATION-END: DEF-03-DOMAIN
% CF-DEFINITION-ENTRY-END: DEF-03-DOMAIN

% CF-DEFINITION-ENTRY-BEGIN: DEF-04-PARTIAL
\subsubsection{\texttt{\detokenize{DEF-04-PARTIAL}}}
\noindent\textit{Formal identifier:} \texttt{\detokenize{slp_complex_partial_derivative}}

\reviewlabel{Isabelle code}
% CF-ISABELLE-BEGIN: DEF-04-PARTIAL
\begin{lstlisting}
definition slp_complex_partial_derivative ::
  "slp_scalar_field \<Rightarrow> 2 \<Rightarrow> slp_point \<Rightarrow> complex"
where
  "slp_complex_partial_derivative phi i x =
    frechet_derivative phi (at x) (axis i 1)"
\end{lstlisting}
% CF-ISABELLE-END: DEF-04-PARTIAL

\reviewlabel{Independent English translation of the Isabelle code}
% CF-BACK-TRANSLATION-BEGIN: DEF-04-PARTIAL
For coordinate $i$, evaluate the Fréchet derivative of the complex field in the direction of the standard axis vector.
% CF-BACK-TRANSLATION-END: DEF-04-PARTIAL
% CF-DEFINITION-ENTRY-END: DEF-04-PARTIAL

% CF-DEFINITION-ENTRY-BEGIN: DEF-05-GRADIENT
\subsubsection{\texttt{\detokenize{DEF-05-GRADIENT}}}
\noindent\textit{Formal identifier:} \texttt{\detokenize{slp_classical_gradient}}

\reviewlabel{Isabelle code}
% CF-ISABELLE-BEGIN: DEF-05-GRADIENT
\begin{lstlisting}
definition slp_classical_gradient ::
  "slp_scalar_field \<Rightarrow> slp_gradient_field"
where
  "slp_classical_gradient phi x =
    (\<chi> i. slp_complex_partial_derivative phi i x)"
\end{lstlisting}
% CF-ISABELLE-END: DEF-05-GRADIENT

\reviewlabel{Independent English translation of the Isabelle code}
% CF-BACK-TRANSLATION-BEGIN: DEF-05-GRADIENT
At every point, assemble the two coordinate Fréchet derivatives into a complex coordinate vector.
% CF-BACK-TRANSLATION-END: DEF-05-GRADIENT
% CF-DEFINITION-ENTRY-END: DEF-05-GRADIENT

% CF-DEFINITION-ENTRY-BEGIN: DEF-06-TEST
\subsubsection{\texttt{\detokenize{DEF-06-TEST}}}
\noindent\textit{Formal identifier:} \texttt{\detokenize{slp_test_function_on}}

\reviewlabel{Isabelle code}
% CF-ISABELLE-BEGIN: DEF-06-TEST
\begin{lstlisting}
definition slp_test_function_on ::
  "slp_point set \<Rightarrow> slp_scalar_field \<Rightarrow> bool"
where
  "slp_test_function_on U phi \<longleftrightarrow>
    smooth_on UNIV phi \<and>
    compact (closure {x. phi x \<noteq> 0}) \<and>
    closure {x. phi x \<noteq> 0} \<subseteq> U"
\end{lstlisting}
% CF-ISABELLE-END: DEF-06-TEST

\reviewlabel{Independent English translation of the Isabelle code}
% CF-BACK-TRANSLATION-BEGIN: DEF-06-TEST
The predicate is the project model of $C_c^\infty(U)$: smoothness holds on the whole plane, while the closure of the set where the function is nonzero is compact and lies inside $U$.
% CF-BACK-TRANSLATION-END: DEF-06-TEST
% CF-DEFINITION-ENTRY-END: DEF-06-TEST

% CF-DEFINITION-ENTRY-BEGIN: DEF-07-WEAK-GRADIENT
\subsubsection{\texttt{\detokenize{DEF-07-WEAK-GRADIENT}}}
\noindent\textit{Formal identifier:} \texttt{\detokenize{slp_weak_gradient_on}}

\reviewlabel{Isabelle code}
% CF-ISABELLE-BEGIN: DEF-07-WEAK-GRADIENT
\begin{lstlisting}
definition slp_weak_gradient_on ::
  "slp_point set \<Rightarrow> slp_scalar_field \<Rightarrow>
    slp_gradient_field \<Rightarrow> bool"
where
  "slp_weak_gradient_on U u Du \<longleftrightarrow>
    (\<forall>phi. slp_test_function_on U phi \<longrightarrow>
      (\<forall>i.
        set_integrable lborel U
          (\<lambda>x. u x * slp_complex_partial_derivative phi i x) \<and>
        set_integrable lborel U (\<lambda>x. Du x $ i * phi x) \<and>
        set_lebesgue_integral lborel U
          (\<lambda>x. u x * slp_complex_partial_derivative phi i x) =
        - set_lebesgue_integral lborel U (\<lambda>x. Du x $ i * phi x)))"
\end{lstlisting}
% CF-ISABELLE-END: DEF-07-WEAK-GRADIENT

\reviewlabel{Independent English translation of the Isabelle code}
% CF-BACK-TRANSLATION-BEGIN: DEF-07-WEAK-GRADIENT
For all compactly supported smooth complex tests and both Cartesian directions, the distributional integration-by-parts identity holds, with explicit integrability guards on both sides.
% CF-BACK-TRANSLATION-END: DEF-07-WEAK-GRADIENT
% CF-DEFINITION-ENTRY-END: DEF-07-WEAK-GRADIENT

% CF-DEFINITION-ENTRY-BEGIN: DEF-08-H1-PAIR
\subsubsection{\texttt{\detokenize{DEF-08-H1-PAIR}}}
\noindent\textit{Formal identifier:} \texttt{\detokenize{slp_h1_pair_on}}

\reviewlabel{Isabelle code}
% CF-ISABELLE-BEGIN: DEF-08-H1-PAIR
\begin{lstlisting}
definition slp_h1_pair_on ::
  "slp_point set \<Rightarrow> slp_scalar_field \<Rightarrow>
    slp_gradient_field \<Rightarrow> bool"
where
  "slp_h1_pair_on U u Du \<longleftrightarrow>
    u \<in> borel_measurable (restrict_space lborel U) \<and>
    Du \<in> borel_measurable (restrict_space lborel U) \<and>
    slp_weak_gradient_on U u Du \<and>
    set_integrable lborel U (\<lambda>x. norm (u x) ^ 2) \<and>
    set_integrable lborel U (\<lambda>x. norm (Du x) ^ 2)"
\end{lstlisting}
% CF-ISABELLE-END: DEF-08-H1-PAIR

\reviewlabel{Independent English translation of the Isabelle code}
% CF-BACK-TRANSLATION-BEGIN: DEF-08-H1-PAIR
This is representative-level complex $H^1$: measurable function and vector field, the recorded weak-gradient identity, and finite $L^2$ mass for the function and gradient.
% CF-BACK-TRANSLATION-END: DEF-08-H1-PAIR
% CF-DEFINITION-ENTRY-END: DEF-08-H1-PAIR

% CF-DEFINITION-ENTRY-BEGIN: DEF-09-H1-DATA
\subsubsection{\texttt{\detokenize{DEF-09-H1-DATA}}}
\noindent\textit{Formal identifier:} \texttt{\detokenize{slp_h1_data_on}}

\reviewlabel{Isabelle code}
% CF-ISABELLE-BEGIN: DEF-09-H1-DATA
\begin{lstlisting}
definition slp_h1_data_on :: "slp_point set \<Rightarrow> slp_h1_data \<Rightarrow> bool"
where
  "slp_h1_data_on U F \<longleftrightarrow>
    slp_h1_pair_on U (fst F) (snd F)"
\end{lstlisting}
% CF-ISABELLE-END: DEF-09-H1-DATA

\reviewlabel{Independent English translation of the Isabelle code}
% CF-BACK-TRANSLATION-BEGIN: DEF-09-H1-DATA
The pair-valued representation is accepted on $U$ precisely when its scalar and gradient components constitute project $H^1$ data there.
% CF-BACK-TRANSLATION-END: DEF-09-H1-DATA
% CF-DEFINITION-ENTRY-END: DEF-09-H1-DATA

% CF-DEFINITION-ENTRY-BEGIN: DEF-10-H1-DISTANCE
\subsubsection{\texttt{\detokenize{DEF-10-H1-DISTANCE}}}
\noindent\textit{Formal identifier:} \texttt{\detokenize{slp_h1_squared_distance}}

\reviewlabel{Isabelle code}
% CF-ISABELLE-BEGIN: DEF-10-H1-DISTANCE
\begin{lstlisting}
definition slp_h1_squared_distance ::
  "slp_point set \<Rightarrow> slp_scalar_field \<Rightarrow>
    slp_gradient_field \<Rightarrow> slp_scalar_field \<Rightarrow>
    slp_gradient_field \<Rightarrow> real"
where
  "slp_h1_squared_distance U u Du v Dv =
    set_lebesgue_integral lborel U
      (\<lambda>x. norm (u x - v x) ^ 2 + norm (Du x - Dv x) ^ 2)"
\end{lstlisting}
% CF-ISABELLE-END: DEF-10-H1-DISTANCE

\reviewlabel{Independent English translation of the Isabelle code}
% CF-BACK-TRANSLATION-BEGIN: DEF-10-H1-DISTANCE
This functional is the graph-norm square of the difference between two scalar/gradient representatives, using the restricted planar Lebesgue integral.
% CF-BACK-TRANSLATION-END: DEF-10-H1-DISTANCE
% CF-DEFINITION-ENTRY-END: DEF-10-H1-DISTANCE

% CF-DEFINITION-ENTRY-BEGIN: DEF-11-H10-PAIR
\subsubsection{\texttt{\detokenize{DEF-11-H10-PAIR}}}
\noindent\textit{Formal identifier:} \texttt{\detokenize{slp_h1_zero_pair_on}}

\reviewlabel{Isabelle code}
% CF-ISABELLE-BEGIN: DEF-11-H10-PAIR
\begin{lstlisting}
definition slp_h1_zero_pair_on ::
  "slp_point set \<Rightarrow> slp_scalar_field \<Rightarrow>
    slp_gradient_field \<Rightarrow> bool"
where
  "slp_h1_zero_pair_on U u Du \<longleftrightarrow>
    slp_h1_pair_on U u Du \<and>
    (\<exists>seq :: nat \<Rightarrow> slp_scalar_field.
      (\<forall>m. slp_test_function_on U (seq m) \<and>
        slp_h1_pair_on U (seq m) (slp_classical_gradient (seq m))) \<and>
      (\<forall>epsilon>0. \<exists>N. \<forall>m\<ge>N.
        slp_h1_squared_distance U (seq m)
          (slp_classical_gradient (seq m)) u Du < epsilon))"
\end{lstlisting}
% CF-ISABELLE-END: DEF-11-H10-PAIR

\reviewlabel{Independent English translation of the Isabelle code}
% CF-BACK-TRANSLATION-BEGIN: DEF-11-H10-PAIR
The zero-boundary Sobolev predicate is the closure of $C_c^\infty(U)$ in the project $H^1$ graph norm, retaining both the scalar representative and its supplied weak gradient.
% CF-BACK-TRANSLATION-END: DEF-11-H10-PAIR
% CF-DEFINITION-ENTRY-END: DEF-11-H10-PAIR

% CF-DEFINITION-ENTRY-BEGIN: DEF-12-H10-DATA
\subsubsection{\texttt{\detokenize{DEF-12-H10-DATA}}}
\noindent\textit{Formal identifier:} \texttt{\detokenize{slp_h1_zero_data_on}}

\reviewlabel{Isabelle code}
% CF-ISABELLE-BEGIN: DEF-12-H10-DATA
\begin{lstlisting}
definition slp_h1_zero_data_on ::
  "slp_point set \<Rightarrow> slp_h1_data \<Rightarrow> bool"
where
  "slp_h1_zero_data_on U F \<longleftrightarrow>
    slp_h1_zero_pair_on U (fst F) (snd F)"
\end{lstlisting}
% CF-ISABELLE-END: DEF-12-H10-DATA

\reviewlabel{Independent English translation of the Isabelle code}
% CF-BACK-TRANSLATION-BEGIN: DEF-12-H10-DATA
The ordered $H^1$ datum is zero-boundary data when its scalar and gradient components belong to the test-function closure just defined.
% CF-BACK-TRANSLATION-END: DEF-12-H10-DATA
% CF-DEFINITION-ENTRY-END: DEF-12-H10-DATA

% CF-DEFINITION-ENTRY-BEGIN: DEF-13-LP
\subsubsection{\texttt{\detokenize{DEF-13-LP}}}
\noindent\textit{Formal identifier:} \texttt{\detokenize{slp_complex_lp_on}}

\reviewlabel{Isabelle code}
% CF-ISABELLE-BEGIN: DEF-13-LP
\begin{lstlisting}
definition slp_complex_lp_on ::
  "real \<Rightarrow> slp_point set \<Rightarrow> slp_scalar_field \<Rightarrow> bool"
where
  "slp_complex_lp_on p U V \<longleftrightarrow>
    V \<in> borel_measurable (restrict_space lborel U) \<and>
    set_integrable lborel U (\<lambda>x. norm (V x) powr p)"
\end{lstlisting}
% CF-ISABELLE-END: DEF-13-LP

\reviewlabel{Independent English translation of the Isabelle code}
% CF-BACK-TRANSLATION-BEGIN: DEF-13-LP
The predicate records a measurable complex representative on $U$ whose $p$th power of the modulus is integrable over $U$.
% CF-BACK-TRANSLATION-END: DEF-13-LP
% CF-DEFINITION-ENTRY-END: DEF-13-LP

% CF-DEFINITION-ENTRY-BEGIN: DEF-14-WEAK-FORM-INTEGRABLE
\subsubsection{\texttt{\detokenize{DEF-14-WEAK-FORM-INTEGRABLE}}}
\noindent\textit{Formal identifier:} \texttt{\detokenize{slp_weak_form_integrable}}

\reviewlabel{Isabelle code}
% CF-ISABELLE-BEGIN: DEF-14-WEAK-FORM-INTEGRABLE
\begin{lstlisting}
definition slp_weak_form_integrable ::
  "slp_point set \<Rightarrow> slp_scalar_field \<Rightarrow>
    slp_h1_data \<Rightarrow> slp_h1_data \<Rightarrow> bool"
where
  "slp_weak_form_integrable U V F G \<longleftrightarrow>
    set_integrable lborel U
      (\<lambda>x. \<Sum>i\<in>UNIV. snd F x $ i * snd G x $ i) \<and>
    set_integrable lborel U
      (\<lambda>x. V x * fst F x * fst G x)"
\end{lstlisting}
% CF-ISABELLE-END: DEF-14-WEAK-FORM-INTEGRABLE

\reviewlabel{Independent English translation of the Isabelle code}
% CF-BACK-TRANSLATION-BEGIN: DEF-14-WEAK-FORM-INTEGRABLE
This guard requires integrability of the sum of coordinatewise gradient products and, separately, of $V(F_0)(G_0)$; neither product uses complex conjugation.
% CF-BACK-TRANSLATION-END: DEF-14-WEAK-FORM-INTEGRABLE
% CF-DEFINITION-ENTRY-END: DEF-14-WEAK-FORM-INTEGRABLE

% CF-DEFINITION-ENTRY-BEGIN: DEF-15-WEAK-FORM
\subsubsection{\texttt{\detokenize{DEF-15-WEAK-FORM}}}
\noindent\textit{Formal identifier:} \texttt{\detokenize{slp_weak_form}}

\reviewlabel{Isabelle code}
% CF-ISABELLE-BEGIN: DEF-15-WEAK-FORM
\begin{lstlisting}
definition slp_weak_form ::
  "slp_point set \<Rightarrow> slp_scalar_field \<Rightarrow>
    slp_h1_data \<Rightarrow> slp_h1_data \<Rightarrow> complex"
where
  "slp_weak_form U V F G =
    set_lebesgue_integral lborel U
      (\<lambda>x. \<Sum>i\<in>UNIV. snd F x $ i * snd G x $ i) +
    set_lebesgue_integral lborel U
      (\<lambda>x. V x * fst F x * fst G x)"
\end{lstlisting}
% CF-ISABELLE-END: DEF-15-WEAK-FORM

\reviewlabel{Independent English translation of the Isabelle code}
% CF-BACK-TRANSLATION-BEGIN: DEF-15-WEAK-FORM
For two data pairs, add the restricted integral of the coordinatewise product of their recorded gradients to the restricted integral of the potential times both scalar components, without conjugation.
% CF-BACK-TRANSLATION-END: DEF-15-WEAK-FORM
% CF-DEFINITION-ENTRY-END: DEF-15-WEAK-FORM

% CF-DEFINITION-ENTRY-BEGIN: DEF-16-WEAK-SOLUTION
\subsubsection{\texttt{\detokenize{DEF-16-WEAK-SOLUTION}}}
\noindent\textit{Formal identifier:} \texttt{\detokenize{slp_weak_solution}}

\reviewlabel{Isabelle code}
% CF-ISABELLE-BEGIN: DEF-16-WEAK-SOLUTION
\begin{lstlisting}
definition slp_weak_solution ::
  "slp_point set \<Rightarrow> slp_scalar_field \<Rightarrow>
    slp_h1_data \<Rightarrow> bool"
where
  "slp_weak_solution U V F \<longleftrightarrow>
    slp_h1_data_on U F \<and>
    (\<forall>G. slp_h1_zero_data_on U G \<longrightarrow>
      slp_weak_form_integrable U V F G \<and>
      slp_weak_form U V F G = 0)"
\end{lstlisting}
% CF-ISABELLE-END: DEF-16-WEAK-SOLUTION

\reviewlabel{Independent English translation of the Isabelle code}
% CF-BACK-TRANSLATION-BEGIN: DEF-16-WEAK-SOLUTION
The predicate expresses $(-\Delta+V)u=0$ in the project complex-bilinear convention: the candidate is $H^1$, and its guarded weak form vanishes for all zero-boundary $H^1$ tests.
% CF-BACK-TRANSLATION-END: DEF-16-WEAK-SOLUTION
% CF-DEFINITION-ENTRY-END: DEF-16-WEAK-SOLUTION

% CF-DEFINITION-ENTRY-BEGIN: DEF-17-AE-EQUALITY
\subsubsection{\texttt{\detokenize{DEF-17-AE-EQUALITY}}}
\noindent\textit{Formal identifier:} \texttt{\detokenize{slp_potential_ae_equal}}

\reviewlabel{Isabelle code}
% CF-ISABELLE-BEGIN: DEF-17-AE-EQUALITY
\begin{lstlisting}
definition slp_potential_ae_equal ::
  "slp_point set \<Rightarrow> slp_scalar_field \<Rightarrow>
    slp_scalar_field \<Rightarrow> bool"
where
  "slp_potential_ae_equal U V V_tilde \<longleftrightarrow>
    (AE x in restrict_space lborel U. V x = V_tilde x)"
\end{lstlisting}
% CF-ISABELLE-END: DEF-17-AE-EQUALITY

\reviewlabel{Independent English translation of the Isabelle code}
% CF-BACK-TRANSLATION-BEGIN: DEF-17-AE-EQUALITY
The conclusion ignores values outside $U$ and null-set changes inside it: $V(x)=\widetilde V(x)$ holds almost everywhere in the restricted measure space.
% CF-BACK-TRANSLATION-END: DEF-17-AE-EQUALITY
% CF-DEFINITION-ENTRY-END: DEF-17-AE-EQUALITY

% CF-DEFINITION-ENTRY-BEGIN: DEF-18-ALESSANDRINI
\subsubsection{\texttt{\detokenize{DEF-18-ALESSANDRINI}}}
\noindent\textit{Formal identifier:} \texttt{\detokenize{slp_alessandrini_orthogonality}}

\reviewlabel{Isabelle code}
% CF-ISABELLE-BEGIN: DEF-18-ALESSANDRINI
\begin{lstlisting}
definition slp_alessandrini_orthogonality ::
  "slp_point set \<Rightarrow> slp_scalar_field \<Rightarrow>
    slp_scalar_field \<Rightarrow> bool"
where
  "slp_alessandrini_orthogonality Omega V V_tilde \<longleftrightarrow>
    (\<forall>F G.
      slp_weak_solution Omega V F \<and>
      slp_weak_solution Omega V_tilde G
      \<longrightarrow>
      set_integrable lborel Omega
        (\<lambda>x. (V x - V_tilde x) * fst F x * fst G x) \<and>
      set_lebesgue_integral lborel Omega
        (\<lambda>x. (V x - V_tilde x) * fst F x * fst G x) = 0)"
\end{lstlisting}
% CF-ISABELLE-END: DEF-18-ALESSANDRINI

\reviewlabel{Independent English translation of the Isabelle code}
% CF-BACK-TRANSLATION-BEGIN: DEF-18-ALESSANDRINI
The predicate universally ranges over the full cross-product of the two weak-solution spaces and requires a genuinely integrable, unconjugated interior pairing with value zero for each pair.
% CF-BACK-TRANSLATION-END: DEF-18-ALESSANDRINI
% CF-DEFINITION-ENTRY-END: DEF-18-ALESSANDRINI
% CF-DEFINITIONS-END

\clearpage
\section{Statements}
\subsection{Advertised conditional conclusion}
% CF-MANUSCRIPT-CLAIMS-BEGIN
% CF-MANUSCRIPT-ENTRY-BEGIN: MAIN-001
\subsubsection{\texttt{\detokenize{MAIN-001}}}
\noindent\textit{Formal identifier:} \texttt{\detokenize{slp_alessandrini_uniqueness_claim}}

\reviewlabel{English original and exact authority locator}
% CF-MANUSCRIPT-ORIGINAL-BEGIN: MAIN-001
\begin{sourcequote}
The frozen user-selected statement specification v002, based on the official
manuscript's weak Alessandrini identity at lines 1115--1152 and equation
\texttt{weak-alessandrini-zero}, states: let $\Omega$ be a bounded smooth
domain in the real plane, let $1<p<\infty$, and let $V,\widetilde V$ be
possibly complex-valued members of $L^p(\Omega)$. Suppose that for every
$H^1(\Omega)$ weak solution $u$ of $(-\Delta+V)u=0$ and every
$H^1(\Omega)$ weak solution $\widetilde u$ of
$(-\Delta+\widetilde V)\widetilde u=0$, the complex-bilinear product
$(V-\widetilde V)u\widetilde u$ is integrable on $\Omega$ and has integral
zero. Then $V=\widetilde V$ almost everywhere in $\Omega$. No complex
conjugate occurs. This is the user-authorized conditional replacement target,
not manuscript Theorem \texttt{main-thm}; non-eigenvalue and DN-map premises
are deliberately outside v002.
\end{sourcequote}
% CF-MANUSCRIPT-ORIGINAL-END: MAIN-001

\reviewlabel{Isabelle code}
% CF-ISABELLE-BEGIN: MAIN-001
\begin{lstlisting}
definition slp_alessandrini_uniqueness_claim ::
  "real \<Rightarrow> slp_point set \<Rightarrow> slp_scalar_field \<Rightarrow>
    slp_scalar_field \<Rightarrow> bool"
where
  "slp_alessandrini_uniqueness_claim p Omega V V_tilde \<longleftrightarrow>
    (1 < p \<and>
     slp_bounded_smooth_domain Omega \<and>
     slp_complex_lp_on p Omega V \<and>
     slp_complex_lp_on p Omega V_tilde \<and>
     slp_alessandrini_orthogonality Omega V V_tilde)
    \<longrightarrow> slp_potential_ae_equal Omega V V_tilde"
\end{lstlisting}
% CF-ISABELLE-END: MAIN-001

\reviewlabel{Independent English translation of the Isabelle code}
% CF-BACK-TRANSLATION-BEGIN: MAIN-001
The Boolean proposition is an implication, universally parameterized by
$p,\Omega,V,\widetilde V$. Its antecedent combines the finite exponent
condition $1<p$, the exact bounded-smooth-domain predicate, both local complex
$L^p$ hypotheses, and universal bilinear orthogonality over the full cross
product of the two weak-solution spaces. Its conclusion is equality only
almost everywhere inside $\Omega$. It assumes neither DN equality nor
non-eigenvalue conditions and says nothing about representatives outside the
domain.
% CF-BACK-TRANSLATION-END: MAIN-001
% CF-MANUSCRIPT-ENTRY-END: MAIN-001
% CF-MANUSCRIPT-CLAIMS-END

\subsection{Cited results}
Each entry first states the complete source material used, then the exact
derived interface result, its Isabelle declaration, and an independently
produced English back-translation.
% CF-LITERATURE-ENTRIES-BEGIN
% CF-LITERATURE-ENTRY-BEGIN: AIM-PLANAR-BEURLING-LP-BOUND
\subsubsection{\texttt{\detokenize{AIM-PLANAR-BEURLING-LP-BOUND}}}
\noindent\textit{Formal identifier:} \texttt{\detokenize{aim_planar_beurling_lp_claim}}.\par
\noindent\textit{Relationship to source:} derived.

\reviewlabel{Source attribution and locator}
% CF-SOURCE-ATTRIBUTION-BEGIN: AIM-PLANAR-BEURLING-LP-BOUND/SOURCE-1
Kari Astala, Tadeusz Iwaniec, and Gaven Martin, Elliptic Partial Differential Equations and Quasiconformal Mappings in the Plane, Princeton Mathematical Series 48, Princeton University Press, 2009, formula (4.17), Theorem 4.0.10, and Theorem 4.5.3 with (4.84), printed pp. 94, 97-98, and 129, ISBN 978-0-691-13777-3
% CF-SOURCE-ATTRIBUTION-END: AIM-PLANAR-BEURLING-LP-BOUND/SOURCE-1

\reviewlabel{Quoted source theorem statement}
% CF-SOURCE-STATEMENT-BEGIN: AIM-PLANAR-BEURLING-LP-BOUND/SOURCE-1
\begin{sourcequote}
AIM formula (4.17) defines the Beurling transform by the principal value $-(1/\pi)\lim_{\varepsilon\downarrow0}\int_{|z-\tau|>\varepsilon} f(\tau)/(z-\tau)^2\,d\tau$. Theorem 4.0.10 states that for $f\in L^p(\mathbb C)$, $1<p<\infty$, this limit exists almost everywhere. Theorem 4.5.3 states that $S$ extends continuously on $L^p(\mathbb C)$ for every $1<p<\infty$ and satisfies $\|Sf\|_p\leq S_p\|f\|_p$.
\end{sourcequote}
% CF-SOURCE-STATEMENT-END: AIM-PLANAR-BEURLING-LP-BOUND/SOURCE-1

\reviewlabel{Formalized derived mathematical result}
% CF-DERIVED-RESULT-BEGIN: AIM-PLANAR-BEURLING-LP-BOUND
Consequently, for each real $p>1$ there is $S_p>0$ such that every measurable complex $L^p$ representative has the displayed principal value almost everywhere; choosing that value where it exists and zero on the exceptional null set produces another $L^p$ representative with norm at most $S_p$ times the original norm.
% CF-DERIVED-RESULT-END: AIM-PLANAR-BEURLING-LP-BOUND

\reviewlabel{Isabelle code}
% CF-ISABELLE-BEGIN: AIM-PLANAR-BEURLING-LP-BOUND
\begin{lstlisting}
definition aim_planar_beurling_truncation ::
  "aim_planar_field \<Rightarrow> real \<Rightarrow> aim_planar_point \<Rightarrow> complex"
where
  "aim_planar_beurling_truncation f epsilon z =
    - inverse (of_real pi :: complex) *
      integral\<^sup>L lborel
        (\<lambda>y. if epsilon <
              norm (aim_point_as_complex z - aim_point_as_complex y)
          then f y * inverse
            ((aim_point_as_complex z - aim_point_as_complex y)\<^sup>2)
          else 0)"

definition aim_planar_beurling_pv_exists ::
  "aim_planar_field \<Rightarrow> aim_planar_point \<Rightarrow> bool"
where
  "aim_planar_beurling_pv_exists f z \<longleftrightarrow>
    (\<exists>w::complex.
      ((\<lambda>epsilon::real.
          aim_planar_beurling_truncation f epsilon z) \<longlongrightarrow> w)
        (at_right 0))"

definition aim_planar_beurling_transform ::
  "aim_planar_field \<Rightarrow> aim_planar_field"
where
  "aim_planar_beurling_transform f z =
    (if aim_planar_beurling_pv_exists f z then
      Lim (at_right 0)
        (\<lambda>epsilon::real.
          aim_planar_beurling_truncation f epsilon z)
    else 0)"

definition aim_planar_beurling_lp_claim :: bool
where
  "aim_planar_beurling_lp_claim \<longleftrightarrow>
    (\<forall>p::real. 1 < p \<longrightarrow>
      (\<exists>S_p::real. 0 < S_p \<and>
        (\<forall>f. aim_complex_lp_on_plane p f \<longrightarrow>
          (AE z in lborel. aim_planar_beurling_pv_exists f z) \<and>
          aim_complex_lp_on_plane p (aim_planar_beurling_transform f) \<and>
          aim_complex_lp_norm p (aim_planar_beurling_transform f)
            \<le> S_p * aim_complex_lp_norm p f)))"
\end{lstlisting}
% CF-ISABELLE-END: AIM-PLANAR-BEURLING-LP-BOUND

\reviewlabel{Independent English translation of the Isabelle code}
% CF-BACK-TRANSLATION-BEGIN: AIM-PLANAR-BEURLING-LP-BOUND
Given $p>1$, a constant depending only on $p$ bounds the formal Beurling transform of every admissible field. The claim also asserts almost-everywhere existence of the positive-radius principal value and measurable $L^p$ membership of the representative that is set to zero off that full-measure set.
% CF-BACK-TRANSLATION-END: AIM-PLANAR-BEURLING-LP-BOUND
% CF-LITERATURE-ENTRY-END: AIM-PLANAR-BEURLING-LP-BOUND

% CF-LITERATURE-ENTRY-BEGIN: AIM-PLANAR-CAUCHY-BEURLING-DERIVATIVES
\subsubsection{\texttt{\detokenize{AIM-PLANAR-CAUCHY-BEURLING-DERIVATIVES}}}
\noindent\textit{Formal identifier:} \texttt{\detokenize{aim_planar_cauchy_beurling_derivatives_claim}}.\par
\noindent\textit{Relationship to source:} derived.

\reviewlabel{Source attribution and locator}
% CF-SOURCE-ATTRIBUTION-BEGIN: AIM-PLANAR-CAUCHY-BEURLING-DERIVATIVES/SOURCE-1
Kari Astala, Tadeusz Iwaniec, and Gaven Martin, Elliptic Partial Differential Equations and Quasiconformal Mappings in the Plane, Princeton Mathematical Series 48, Princeton University Press, 2009, Definition 4.0.9 and formulas (4.6), (4.8), (4.16)-(4.17), Theorem 4.0.10, the Lp extension discussion on printed pp. 96-97, and Theorem 4.5.3, printed pp. 93-98 and 129, ISBN 978-0-691-13777-3
% CF-SOURCE-ATTRIBUTION-END: AIM-PLANAR-CAUCHY-BEURLING-DERIVATIVES/SOURCE-1

\reviewlabel{Quoted source theorem statement}
% CF-SOURCE-STATEMENT-BEGIN: AIM-PLANAR-CAUCHY-BEURLING-DERIVATIVES/SOURCE-1
\begin{sourcequote}
AIM Definition 4.0.9 and (4.6) give $Cf(z)=\pi^{-1}\int f(\tau)/(z-\tau)\,d\tau$; (4.8) gives $\bar\partial C=I$ on test functions; (4.16)--(4.17) give $\partial(Cf)=Sf$ with the principal-value Beurling normalization. The discussion after (4.25), together with Theorems 4.0.10 and 4.5.3, extends these identities to $L^p(\mathbb C)$, $1<p<\infty$; the convergence discussion also makes $Cf$ locally integrable almost everywhere for compactly supported integrable $f$.
\end{sourcequote}
% CF-SOURCE-STATEMENT-END: AIM-PLANAR-CAUCHY-BEURLING-DERIVATIVES/SOURCE-1

\reviewlabel{Formalized derived mathematical result}
% CF-DERIVED-RESULT-BEGIN: AIM-PLANAR-CAUCHY-BEURLING-DERIVATIVES
For a bounded-support measurable $f\in L^p(\mathbb R^2)$, $p>1$, the explicit Cauchy representative has distributional Cartesian gradient $(f+Sf,\, i(Sf-f))$. This is precisely the simultaneous pair of Wirtinger identities $\bar\partial(Cf)=f$ and $\partial(Cf)=Sf$, with every test-product integrability condition made explicit.
% CF-DERIVED-RESULT-END: AIM-PLANAR-CAUCHY-BEURLING-DERIVATIVES

\reviewlabel{Isabelle code}
% CF-ISABELLE-BEGIN: AIM-PLANAR-CAUCHY-BEURLING-DERIVATIVES
\begin{lstlisting}
definition aim_planar_cauchy_beurling_gradient ::
  "aim_planar_field \<Rightarrow> slp_gradient_field"
where
  "aim_planar_cauchy_beurling_gradient f x =
    (\<chi> i. if i = 0 then
      f x + aim_planar_beurling_transform f x
    else
      \<i> * (aim_planar_beurling_transform f x - f x))"

definition aim_planar_cauchy_beurling_derivatives_claim :: bool
where
  "aim_planar_cauchy_beurling_derivatives_claim \<longleftrightarrow>
    (\<forall>p::real. 1 < p \<longrightarrow>
      (\<forall>f. aim_complex_lp_on_plane p f \<and> bounded {x. f x \<noteq> 0}
        \<longrightarrow>
        slp_weak_gradient_on UNIV
          (aim_planar_cauchy_transform f)
          (aim_planar_cauchy_beurling_gradient f)))"
\end{lstlisting}
% CF-ISABELLE-END: AIM-PLANAR-CAUCHY-BEURLING-DERIVATIVES

\reviewlabel{Independent English translation of the Isabelle code}
% CF-BACK-TRANSLATION-BEGIN: AIM-PLANAR-CAUCHY-BEURLING-DERIVATIVES
The claim says that the normalized Cauchy transform of any bounded-support $L^p$ input has, in distributions, first Cartesian derivative $f+Sf$ and second derivative $i(Sf-f)$. The conclusion is guarded weak-gradient equality on all of $\mathbb R^2$.
% CF-BACK-TRANSLATION-END: AIM-PLANAR-CAUCHY-BEURLING-DERIVATIVES
% CF-LITERATURE-ENTRY-END: AIM-PLANAR-CAUCHY-BEURLING-DERIVATIVES

% CF-LITERATURE-ENTRY-BEGIN: AIM-PLANAR-CAUCHY-TEST-LEFT-INVERSE
\subsubsection{\texttt{\detokenize{AIM-PLANAR-CAUCHY-TEST-LEFT-INVERSE}}}
\noindent\textit{Formal identifier:} \texttt{\detokenize{aim_planar_cauchy_test_left_inverse_claim}}.\par
\noindent\textit{Relationship to source:} derived.

\reviewlabel{Source attribution and locator}
% CF-SOURCE-ATTRIBUTION-BEGIN: AIM-PLANAR-CAUCHY-TEST-LEFT-INVERSE/SOURCE-1
Kari Astala, Tadeusz Iwaniec, and Gaven Martin, Elliptic Partial Differential Equations and Quasiconformal Mappings in the Plane, Princeton Mathematical Series 48, Princeton University Press, 2009, Definition 4.0.9 and formulas (4.6)-(4.8), printed p. 93 / PDF p. 112, ISBN 978-0-691-13777-3
% CF-SOURCE-ATTRIBUTION-END: AIM-PLANAR-CAUCHY-TEST-LEFT-INVERSE/SOURCE-1

\reviewlabel{Quoted source theorem statement}
% CF-SOURCE-STATEMENT-BEGIN: AIM-PLANAR-CAUCHY-TEST-LEFT-INVERSE/SOURCE-1
\begin{sourcequote}
For $\phi\in C_c^\infty(\mathbb C)$, AIM Definition 4.0.9 and (4.6) define $C\phi(z)=\pi^{-1}\int_{\mathbb C}\phi(\tau)/(z-\tau)\,d\tau$. Formula (4.8) states the two-sided identity $C\circ\bar\partial=\bar\partial\circ C=I$ on this test-function class.
\end{sourcequote}
% CF-SOURCE-STATEMENT-END: AIM-PLANAR-CAUCHY-TEST-LEFT-INVERSE/SOURCE-1

\reviewlabel{Formalized derived mathematical result}
% CF-DERIVED-RESULT-BEGIN: AIM-PLANAR-CAUCHY-TEST-LEFT-INVERSE
Taking only the left-inverse half and writing $\bar\partial\phi=(\partial_x\phi+i\partial_y\phi)/2$ in the project coordinates gives $C(\bar\partial\phi)=\phi$ as equality of complex-valued functions on the whole plane.
% CF-DERIVED-RESULT-END: AIM-PLANAR-CAUCHY-TEST-LEFT-INVERSE

\reviewlabel{Isabelle code}
% CF-ISABELLE-BEGIN: AIM-PLANAR-CAUCHY-TEST-LEFT-INVERSE
\begin{lstlisting}
definition aim_planar_classical_dbar ::
  "slp_scalar_field \<Rightarrow> slp_scalar_field"
where
  "aim_planar_classical_dbar phi x =
    (slp_complex_partial_derivative phi 0 x +
      \<i> * slp_complex_partial_derivative phi 1 x) / 2"

definition aim_planar_cauchy_test_left_inverse_claim :: bool
where
  "aim_planar_cauchy_test_left_inverse_claim \<longleftrightarrow>
    (\<forall>phi. slp_test_function_on UNIV phi \<longrightarrow>
      aim_planar_cauchy_transform (aim_planar_classical_dbar phi) = phi)"
\end{lstlisting}
% CF-ISABELLE-END: AIM-PLANAR-CAUCHY-TEST-LEFT-INVERSE

\reviewlabel{Independent English translation of the Isabelle code}
% CF-BACK-TRANSLATION-BEGIN: AIM-PLANAR-CAUCHY-TEST-LEFT-INVERSE
For all project test functions on $\mathbb R^2$, the operator $C$ is a left inverse of the explicitly defined Wirtinger operator $(D_0+iD_1)/2$; the conclusion is extensional equality everywhere.
% CF-BACK-TRANSLATION-END: AIM-PLANAR-CAUCHY-TEST-LEFT-INVERSE
% CF-LITERATURE-ENTRY-END: AIM-PLANAR-CAUCHY-TEST-LEFT-INVERSE

% CF-LITERATURE-ENTRY-BEGIN: AIM-PLANAR-HARDY-LITTLEWOOD-SOBOLEV
\subsubsection{\texttt{\detokenize{AIM-PLANAR-HARDY-LITTLEWOOD-SOBOLEV}}}
\noindent\textit{Formal identifier:} \texttt{\detokenize{aim_planar_hls_cauchy_claim}}.\par
\noindent\textit{Relationship to source:} derived.

\reviewlabel{Source attribution and locator}
% CF-SOURCE-ATTRIBUTION-BEGIN: AIM-PLANAR-HARDY-LITTLEWOOD-SOBOLEV/SOURCE-1
Kari Astala, Tadeusz Iwaniec, and Gaven Martin, Elliptic Partial Differential Equations and Quasiconformal Mappings in the Plane, Princeton Mathematical Series 48, Princeton University Press, 2009, Theorem 4.3.8, printed p. 111; normalization and HLS estimate printed pp. 109-110
% CF-SOURCE-ATTRIBUTION-END: AIM-PLANAR-HARDY-LITTLEWOOD-SOBOLEV/SOURCE-1

\reviewlabel{Quoted source theorem statement}
% CF-SOURCE-STATEMENT-BEGIN: AIM-PLANAR-HARDY-LITTLEWOOD-SOBOLEV/SOURCE-1
\begin{sourcequote}
AIM Theorem 4.3.8 states: for every $1<p<2$, the Cauchy transform is continuous from $L^p(\mathbb C)$ to $L^{2p/(2-p)}(\mathbb C)$; moreover $\|C\phi\|_{2p/(2-p)}\leq C\,\|\phi\|_p/((p-1)(2-p))$ for one absolute constant $C$. The transform uses the normalization $C\phi(z)=\pi^{-1}\int\phi(\tau)/(z-\tau)\,d\tau$.
\end{sourcequote}
% CF-SOURCE-STATEMENT-END: AIM-PLANAR-HARDY-LITTLEWOOD-SOBOLEV/SOURCE-1

\reviewlabel{Formalized derived mathematical result}
% CF-DERIVED-RESULT-BEGIN: AIM-PLANAR-HARDY-LITTLEWOOD-SOBOLEV
At representative level, there is one positive real $C$, independent of $p$ and $f$, such that each measurable complex $L^p$ field with $1<p<2$ is sent by the normalized totalized Cauchy representative to the explicit $L^{2p/(2-p)}$ model with exactly the displayed norm bound.
% CF-DERIVED-RESULT-END: AIM-PLANAR-HARDY-LITTLEWOOD-SOBOLEV

\reviewlabel{Isabelle code}
% CF-ISABELLE-BEGIN: AIM-PLANAR-HARDY-LITTLEWOOD-SOBOLEV
\begin{lstlisting}
type_synonym aim_planar_point = "real ^ 2"
type_synonym aim_planar_field = "aim_planar_point \<Rightarrow> complex"

definition aim_point_as_complex :: "aim_planar_point \<Rightarrow> complex" where
  "aim_point_as_complex x = Complex (x $ 0) (x $ 1)"

definition aim_complex_lp_on_plane ::
  "real \<Rightarrow> aim_planar_field \<Rightarrow> bool"
where
  "aim_complex_lp_on_plane p f \<longleftrightarrow>
    f \<in> borel_measurable lborel \<and>
    integrable lborel (\<lambda>x. norm (f x) powr p)"

definition aim_complex_lp_norm ::
  "real \<Rightarrow> aim_planar_field \<Rightarrow> real"
where
  "aim_complex_lp_norm p f =
    (integral\<^sup>L lborel (\<lambda>x. norm (f x) powr p)) powr (1 / p)"

definition aim_hls_target_exponent :: "real \<Rightarrow> real" where
  "aim_hls_target_exponent p = 2 * p / (2 - p)"

definition aim_planar_cauchy_integrand ::
  "aim_planar_field \<Rightarrow> aim_planar_point \<Rightarrow>
    aim_planar_point \<Rightarrow> complex"
where
  "aim_planar_cauchy_integrand f z y =
    f y * inverse (aim_point_as_complex z - aim_point_as_complex y)"

definition aim_planar_cauchy_transform ::
  "aim_planar_field \<Rightarrow> aim_planar_field"
where
  "aim_planar_cauchy_transform f z =
    inverse (of_real pi :: complex) *
      integral\<^sup>L lborel (aim_planar_cauchy_integrand f z)"

definition aim_planar_hls_cauchy_claim :: bool where
  "aim_planar_hls_cauchy_claim \<longleftrightarrow>
    (\<exists>C::real. 0 < C \<and>
      (\<forall>p f. 1 < p \<and> p < 2 \<and> aim_complex_lp_on_plane p f
        \<longrightarrow>
        aim_complex_lp_on_plane (aim_hls_target_exponent p)
          (aim_planar_cauchy_transform f) \<and>
        aim_complex_lp_norm (aim_hls_target_exponent p)
            (aim_planar_cauchy_transform f)
          \<le> C / ((p - 1) * (2 - p)) * aim_complex_lp_norm p f))"
\end{lstlisting}
% CF-ISABELLE-END: AIM-PLANAR-HARDY-LITTLEWOOD-SOBOLEV

\reviewlabel{Independent English translation of the Isabelle code}
% CF-BACK-TRANSLATION-BEGIN: AIM-PLANAR-HARDY-LITTLEWOOD-SOBOLEV
The formal HLS claim existentially fixes a single constant and then universally quantifies $p$ and $f$. Under $1<p<2$ and plane-wide complex $L^p$ membership, it gives target-space membership for $Cf$ at $q=2p/(2-p)$ together with the source estimate.
% CF-BACK-TRANSLATION-END: AIM-PLANAR-HARDY-LITTLEWOOD-SOBOLEV
% CF-LITERATURE-ENTRY-END: AIM-PLANAR-HARDY-LITTLEWOOD-SOBOLEV

% CF-LITERATURE-ENTRY-BEGIN: AIM-PLANAR-RIESZ-POTENTIAL-HLS
\subsubsection{\texttt{\detokenize{AIM-PLANAR-RIESZ-POTENTIAL-HLS}}}
\noindent\textit{Formal identifier:} \texttt{\detokenize{aim_planar_riesz_hls_claim}}.\par
\noindent\textit{Relationship to source:} derived.

\reviewlabel{Source attribution and locator}
% CF-SOURCE-ATTRIBUTION-BEGIN: AIM-PLANAR-RIESZ-POTENTIAL-HLS/SOURCE-1
Kari Astala, Tadeusz Iwaniec, and Gaven Martin, Elliptic Partial Differential Equations and Quasiconformal Mappings in the Plane, Princeton Mathematical Series 48, Princeton University Press, 2009, displayed Hardy-Littlewood-Sobolev inequality on printed p. 110 / PDF p. 129, ISBN 978-0-691-13777-3
% CF-SOURCE-ATTRIBUTION-END: AIM-PLANAR-RIESZ-POTENTIAL-HLS/SOURCE-1

\reviewlabel{Quoted source theorem statement}
% CF-SOURCE-STATEMENT-BEGIN: AIM-PLANAR-RIESZ-POTENTIAL-HLS/SOURCE-1
\begin{sourcequote}
The Hardy--Littlewood--Sobolev inequality displayed on AIM printed page 110 states that for $1<p<2$ there is an absolute $C$ such that $\|z\mapsto\int_{\mathbb C}|\phi(\tau)|/|z-\tau|\,d\tau\|_{2p/(2-p)}\leq C\|\phi\|_p/((p-1)(2-p))$.
\end{sourcequote}
% CF-SOURCE-STATEMENT-END: AIM-PLANAR-RIESZ-POTENTIAL-HLS/SOURCE-1

\reviewlabel{Formalized derived mathematical result}
% CF-DERIVED-RESULT-BEGIN: AIM-PLANAR-RIESZ-POTENTIAL-HLS
For $1<p<2$ and every explicit complex $L^p$ field, the positive kernel is ordinarily integrable for almost every output point. Its totalized planar potential, with no $1/\pi$ factor, is a measurable real field in $L^{2p/(2-p)}$ and obeys the displayed bound with one absolute positive constant.
% CF-DERIVED-RESULT-END: AIM-PLANAR-RIESZ-POTENTIAL-HLS

\reviewlabel{Isabelle code}
% CF-ISABELLE-BEGIN: AIM-PLANAR-RIESZ-POTENTIAL-HLS
\begin{lstlisting}
definition aim_real_lp_on_plane ::
  "real \<Rightarrow> (aim_planar_point \<Rightarrow> real) \<Rightarrow> bool"
where
  "aim_real_lp_on_plane p g \<longleftrightarrow>
    g \<in> borel_measurable lborel \<and>
    integrable lborel (\<lambda>x. abs (g x) powr p)"

definition aim_real_lp_norm ::
  "real \<Rightarrow> (aim_planar_point \<Rightarrow> real) \<Rightarrow> real"
where
  "aim_real_lp_norm p g =
    (integral\<^sup>L lborel (\<lambda>x. abs (g x) powr p)) powr (1 / p)"

definition aim_planar_riesz_integrand ::
  "aim_planar_field \<Rightarrow> aim_planar_point \<Rightarrow>
    aim_planar_point \<Rightarrow> real"
where
  "aim_planar_riesz_integrand f z y =
    norm (f y) *
      inverse (norm (aim_point_as_complex z - aim_point_as_complex y))"

definition aim_planar_riesz_potential ::
  "aim_planar_field \<Rightarrow> aim_planar_point \<Rightarrow> real"
where
  "aim_planar_riesz_potential f z =
    integral\<^sup>L lborel (aim_planar_riesz_integrand f z)"

definition aim_planar_riesz_hls_claim :: bool where
  "aim_planar_riesz_hls_claim \<longleftrightarrow>
    (\<exists>C::real. 0 < C \<and>
      (\<forall>p f. 1 < p \<and> p < 2 \<and> aim_complex_lp_on_plane p f
        \<longrightarrow>
        (AE z in lborel. integrable lborel (aim_planar_riesz_integrand f z)) \<and>
        aim_real_lp_on_plane (aim_hls_target_exponent p)
          (aim_planar_riesz_potential f) \<and>
        aim_real_lp_norm (aim_hls_target_exponent p)
            (aim_planar_riesz_potential f)
          \<le> C / ((p - 1) * (2 - p)) * aim_complex_lp_norm p f))"
\end{lstlisting}
% CF-ISABELLE-END: AIM-PLANAR-RIESZ-POTENTIAL-HLS

\reviewlabel{Independent English translation of the Isabelle code}
% CF-BACK-TRANSLATION-BEGIN: AIM-PLANAR-RIESZ-POTENTIAL-HLS
The assertion combines almost-everywhere integrability of $|f(y)|/|z-y|$ with measurable target-space membership and a uniform norm bound. It uses the unnormalised positive potential and the exponent $2p/(2-p)$.
% CF-BACK-TRANSLATION-END: AIM-PLANAR-RIESZ-POTENTIAL-HLS
% CF-LITERATURE-ENTRY-END: AIM-PLANAR-RIESZ-POTENTIAL-HLS

% CF-LITERATURE-ENTRY-BEGIN: EVANS-COMPACT-SMOOTH-L1-DENSITY
\subsubsection{\texttt{\detokenize{EVANS-COMPACT-SMOOTH-L1-DENSITY}}}
\noindent\textit{Formal identifier:} \texttt{\detokenize{evans_compact_smooth_l1_density_claim}}.\par
\noindent\textit{Relationship to source:} derived.

\reviewlabel{Source attribution and locator}
% CF-SOURCE-ATTRIBUTION-BEGIN: EVANS-COMPACT-SMOOTH-L1-DENSITY/SOURCE-1
Lawrence C. Evans, Partial Differential Equations, second edition, Graduate Studies in Mathematics 19, American Mathematical Society, 2010, Appendix C.5, standard mollifier printed p. 713 and Theorem 7(i),(iv), printed pp. 714-716
% CF-SOURCE-ATTRIBUTION-END: EVANS-COMPACT-SMOOTH-L1-DENSITY/SOURCE-1

\reviewlabel{Quoted source theorem statement}
% CF-SOURCE-STATEMENT-BEGIN: EVANS-COMPACT-SMOOTH-L1-DENSITY/SOURCE-1
\begin{sourcequote}
Evans Appendix C.5 defines the standard compactly supported smooth mollifier and, in Theorem 7(i),(iv), states that mollifications are smooth on the interior where they are defined and converge to $f$ in $L^p_{\mathrm{loc}}(U)$ for every $1\leq p<\infty$. The mollifier has support in the ball of its scale.
\end{sourcequote}
% CF-SOURCE-STATEMENT-END: EVANS-COMPACT-SMOOTH-L1-DENSITY/SOURCE-1

\reviewlabel{Formalized derived mathematical result}
% CF-DERIVED-RESULT-BEGIN: EVANS-COMPACT-SMOOTH-L1-DENSITY
Truncate a complex $L^1(\mathbb R^n)$ function to a large ball, mollify its real and imaginary parts at one common small scale, and use support propagation plus local-to-global convergence. Thus every complex $L^1$ field is approximated in global $L^1$ norm by a globally smooth compactly supported integrable amplitude.
% CF-DERIVED-RESULT-END: EVANS-COMPACT-SMOOTH-L1-DENSITY

\reviewlabel{Isabelle code}
% CF-ISABELLE-BEGIN: EVANS-COMPACT-SMOOTH-L1-DENSITY
\begin{lstlisting}
definition evans_compact_smooth_l1_density_claim ::
  "'n::finite itself \<Rightarrow> bool"
where
  "evans_compact_smooth_l1_density_claim dimension_type \<longleftrightarrow>
    (\<forall>(F :: real^'n \<Rightarrow> complex).
      integrable lborel F \<longrightarrow>
      (\<forall>epsilon > 0.
        \<exists>u.
          hormander_compact_smooth_amplitude u \<and>
          integral\<^sup>L lborel (\<lambda>x. norm (F x - u x)) < epsilon))"
\end{lstlisting}
% CF-ISABELLE-END: EVANS-COMPACT-SMOOTH-L1-DENSITY

\reviewlabel{Independent English translation of the Isabelle code}
% CF-BACK-TRANSLATION-BEGIN: EVANS-COMPACT-SMOOTH-L1-DENSITY
For any finite coordinate type, compactly supported smooth complex amplitudes are dense in Bochner $L^1$ on Euclidean space, with the approximation expressed by the integral of the pointwise modulus.
% CF-BACK-TRANSLATION-END: EVANS-COMPACT-SMOOTH-L1-DENSITY
% CF-LITERATURE-ENTRY-END: EVANS-COMPACT-SMOOTH-L1-DENSITY

% CF-LITERATURE-ENTRY-BEGIN: EVANS-COMPACT-SMOOTH-LP-DENSITY
\subsubsection{\texttt{\detokenize{EVANS-COMPACT-SMOOTH-LP-DENSITY}}}
\noindent\textit{Formal identifier:} \texttt{\detokenize{evans_compact_smooth_lp_density_claim}}.\par
\noindent\textit{Relationship to source:} derived.

\reviewlabel{Source attribution and locator}
% CF-SOURCE-ATTRIBUTION-BEGIN: EVANS-COMPACT-SMOOTH-LP-DENSITY/SOURCE-1
Lawrence C. Evans, Partial Differential Equations, 2nd ed., Graduate Studies in Mathematics 19, American Mathematical Society, 2010, Appendix C.5 standard mollifier and Theorem 7(i),(iv), printed pp. 713-716, ISBN 978-0-8218-4974-3
% CF-SOURCE-ATTRIBUTION-END: EVANS-COMPACT-SMOOTH-LP-DENSITY/SOURCE-1

\reviewlabel{Quoted source theorem statement}
% CF-SOURCE-STATEMENT-BEGIN: EVANS-COMPACT-SMOOTH-LP-DENSITY/SOURCE-1
\begin{sourcequote}
Evans Appendix C.5 gives a compactly supported smooth standard mollifier and Theorem 7(i),(iv): mollification is smooth in the interior and, for every finite $1\leq p<\infty$, converges to the original function in local $L^p$. The scaled mollifier is supported in the corresponding small ball.
\end{sourcequote}
% CF-SOURCE-STATEMENT-END: EVANS-COMPACT-SMOOTH-LP-DENSITY/SOURCE-1

\reviewlabel{Formalized derived mathematical result}
% CF-DERIVED-RESULT-BEGIN: EVANS-COMPACT-SMOOTH-LP-DENSITY
After large-ball truncation and one common-scale mollification of real and imaginary parts, every measurable complex field with finite $q$th-power integral, $1\leq q<\infty$, has a globally smooth compactly supported approximant with global rooted $L^q$ error below any prescribed positive tolerance.
% CF-DERIVED-RESULT-END: EVANS-COMPACT-SMOOTH-LP-DENSITY

\reviewlabel{Isabelle code}
% CF-ISABELLE-BEGIN: EVANS-COMPACT-SMOOTH-LP-DENSITY
\begin{lstlisting}
definition evans_compact_smooth_lp_density_claim ::
  "'n::finite itself \<Rightarrow> bool"
where
  "evans_compact_smooth_lp_density_claim dimension_type \<longleftrightarrow>
    (\<forall>q (F :: real^'n \<Rightarrow> complex) epsilon.
      1 \<le> q \<and>
      F \<in> borel_measurable lborel \<and>
      integrable lborel (\<lambda>x. norm (F x) powr q) \<and>
      0 < epsilon
      \<longrightarrow>
      (\<exists>u.
        hormander_compact_smooth_amplitude u \<and>
        (integral\<^sup>L lborel
          (\<lambda>x. norm (F x - u x) powr q)) powr (1 / q) < epsilon))"
\end{lstlisting}
% CF-ISABELLE-END: EVANS-COMPACT-SMOOTH-LP-DENSITY

\reviewlabel{Independent English translation of the Isabelle code}
% CF-BACK-TRANSLATION-BEGIN: EVANS-COMPACT-SMOOTH-LP-DENSITY
The claim is finite-dimensional complex $C_c^\infty$ density in $L^q$: it universally quantifies the exponent, field and tolerance, and bounds the $q$th-root of the integrated $q$th power of the approximation error.
% CF-BACK-TRANSLATION-END: EVANS-COMPACT-SMOOTH-LP-DENSITY
% CF-LITERATURE-ENTRY-END: EVANS-COMPACT-SMOOTH-LP-DENSITY

% CF-LITERATURE-ENTRY-BEGIN: EVANS-COMPACT-SUPPORT-W1P-ZERO-DENSITY
\subsubsection{\texttt{\detokenize{EVANS-COMPACT-SUPPORT-W1P-ZERO-DENSITY}}}
\noindent\textit{Formal identifier:} \texttt{\detokenize{evans_compact_support_w1p_zero_density_claim}}.\par
\noindent\textit{Relationship to source:} derived.

\reviewlabel{Source attribution and locator}
% CF-SOURCE-ATTRIBUTION-BEGIN: EVANS-COMPACT-SUPPORT-W1P-ZERO-DENSITY/SOURCE-1
Lawrence C. Evans, Partial Differential Equations, second edition, Graduate Studies in Mathematics 19, American Mathematical Society, 2010, definition of C\_c\textasciicircum{}infinity(U), printed p. 256; definition of W\_0\textasciicircum{}\{k,p\}(U), printed p. 259; Section 5.3.1, Theorem 1 and proof, printed pp. 264-265; Appendix C.5, standard mollifier and Theorem 7(i),(iv), printed pp. 713-716; ISBN 978-0-8218-4974-3
% CF-SOURCE-ATTRIBUTION-END: EVANS-COMPACT-SUPPORT-W1P-ZERO-DENSITY/SOURCE-1

\reviewlabel{Quoted source theorem statement}
% CF-SOURCE-STATEMENT-BEGIN: EVANS-COMPACT-SUPPORT-W1P-ZERO-DENSITY/SOURCE-1
\begin{sourcequote}
For an open Euclidean domain $U\subseteq\mathbb R^n$ and a finite exponent $1\leq p<\infty$, Evans defines $C_c^\infty(U)$ as smooth functions compactly supported in $U$ and $W_0^{k,p}(U)$ as its closure in $W^{k,p}(U)$. Section 5.3.1, Theorem 1 constructs local mollifications of a $W^{k,p}$ function, identifies their derivatives with mollified weak derivatives, and obtains componentwise local $L^p$ convergence; Appendix C.5 supplies the smooth compactly supported mollifier and its finite-$p$ convergence.
\end{sourcequote}
% CF-SOURCE-STATEMENT-END: EVANS-COMPACT-SUPPORT-W1P-ZERO-DENSITY/SOURCE-1

\reviewlabel{Formalized derived mathematical result}
% CF-DERIVED-RESULT-BEGIN: EVANS-COMPACT-SUPPORT-W1P-ZERO-DENSITY
For $p\geq1$ in the plane, if a project $W^{1,p}$ function and both recorded weak-gradient components vanish outside one compact $K\subseteq X$ with $X$ open, sufficiently small mollifications remain compactly supported in $X$ and converge simultaneously in the project graph norm. Hence the pair belongs to project $W^{1,p}_0(X)$.
% CF-DERIVED-RESULT-END: EVANS-COMPACT-SUPPORT-W1P-ZERO-DENSITY

\reviewlabel{Isabelle code}
% CF-ISABELLE-BEGIN: EVANS-COMPACT-SUPPORT-W1P-ZERO-DENSITY
\begin{lstlisting}
definition evans_compact_support_w1p_zero_density_claim :: bool
where
  "evans_compact_support_w1p_zero_density_claim \<longleftrightarrow>
    (\<forall>p X u Du K.
      1 \<le> p \<and>
      open X \<and>
      compact K \<and>
      K \<subseteq> X \<and>
      slp_w1p_pair_on p X u Du \<and>
      (\<forall>x \<in> X - K. u x = 0 \<and> (\<forall>i. Du x $ i = 0))
      \<longrightarrow> slp_w1p_zero_pair_on p X u Du)"
\end{lstlisting}
% CF-ISABELLE-END: EVANS-COMPACT-SUPPORT-W1P-ZERO-DENSITY

\reviewlabel{Independent English translation of the Isabelle code}
% CF-BACK-TRANSLATION-BEGIN: EVANS-COMPACT-SUPPORT-W1P-ZERO-DENSITY
Given an open $X$, compact $K\subseteq X$, and a project weak $W^{1,p}$ pair that vanishes componentwise on $X\setminus K$, the formal smooth-test-function closure predicate follows for all finite real $p\geq1$.
% CF-BACK-TRANSLATION-END: EVANS-COMPACT-SUPPORT-W1P-ZERO-DENSITY
% CF-LITERATURE-ENTRY-END: EVANS-COMPACT-SUPPORT-W1P-ZERO-DENSITY

% CF-LITERATURE-ENTRY-BEGIN: HORMANDER-EUCLIDEAN-L2-FOURIER-PLANCHEREL
\subsubsection{\texttt{\detokenize{HORMANDER-EUCLIDEAN-L2-FOURIER-PLANCHEREL}}}
\noindent\textit{Formal identifier:} \texttt{\detokenize{hormander_euclidean_l2_fourier_plancherel_claim}}.\par
\noindent\textit{Relationship to source:} derived.

\reviewlabel{Source attribution and locator}
% CF-SOURCE-ATTRIBUTION-BEGIN: HORMANDER-EUCLIDEAN-L2-FOURIER-PLANCHEREL/SOURCE-1
Lars Hormander, The Analysis of Linear Partial Differential Operators I: Distribution Theory and Fourier Analysis, first edition, Springer-Verlag, 1983, Definition 7.1.1, formulas (7.1.3) and (7.1.8), Definition 7.1.9, and Theorems 7.1.10-7.1.11, ISBN 3-540-12104-8 / 0-387-12104-8, DOI 10.1007/978-3-642-96750-4.
% CF-SOURCE-ATTRIBUTION-END: HORMANDER-EUCLIDEAN-L2-FOURIER-PLANCHEREL/SOURCE-1

\reviewlabel{Quoted source theorem statement}
% CF-SOURCE-STATEMENT-BEGIN: HORMANDER-EUCLIDEAN-L2-FOURIER-PLANCHEREL/SOURCE-1
\begin{sourcequote}
H\"ormander Definition 7.1.1 and (7.1.3) define $\widehat f(\xi)=\int e^{-i\langle x,\xi\rangle}f(x)\,dx$ for $f\in L^1(\mathbb R^n)$. Parseval (7.1.8) has factor $(2\pi)^{-n}$. Definition 7.1.9 and Theorem 7.1.10 extend the transform to tempered distributions and give $\widehat{\widehat u}(x)=(2\pi)^n u(-x)$. Theorem 7.1.11 states that $u\in L^2$ implies $\widehat u\in L^2$ and that Parseval holds for all $L^2$ functions.
\end{sourcequote}
% CF-SOURCE-STATEMENT-END: HORMANDER-EUCLIDEAN-L2-FOURIER-PLANCHEREL/SOURCE-1

\reviewlabel{Formalized derived mathematical result}
% CF-DERIVED-RESULT-BEGIN: HORMANDER-EUCLIDEAN-L2-FOURIER-PLANCHEREL
Specializing to $n=2$ and selecting measurable representatives gives an operator on explicit complex fields that maps $L^2$ to $L^2$, agrees almost everywhere with the unnormalised negative-exponential integral on $L^1\cap L^2$, scales squared norm by $(2\pi)^2$, and squares to $(2\pi)^2$ times reflection almost everywhere.
% CF-DERIVED-RESULT-END: HORMANDER-EUCLIDEAN-L2-FOURIER-PLANCHEREL

\reviewlabel{Isabelle code}
% CF-ISABELLE-BEGIN: HORMANDER-EUCLIDEAN-L2-FOURIER-PLANCHEREL
\begin{lstlisting}
type_synonym hormander_fourier_point = "real ^ 2"
type_synonym hormander_fourier_field =
  "hormander_fourier_point \<Rightarrow> complex"

definition hormander_fourier_l2_on_plane ::
  "hormander_fourier_field \<Rightarrow> bool"
where
  "hormander_fourier_l2_on_plane f \<longleftrightarrow>
    f \<in> borel_measurable lborel \<and>
    integrable lborel (\<lambda>x. norm (f x) ^ 2)"

definition hormander_fourier_phase ::
  "hormander_fourier_point \<Rightarrow> hormander_fourier_point \<Rightarrow>
    complex"
where
  "hormander_fourier_phase xi x =
    exp (\<i> * of_real (- inner x xi))"

definition hormander_fourier_l1_transform ::
  "hormander_fourier_field \<Rightarrow> hormander_fourier_field"
where
  "hormander_fourier_l1_transform f xi =
    integral\<^sup>L lborel (\<lambda>x. hormander_fourier_phase xi x * f x)"

definition hormander_euclidean_l2_fourier_plancherel_claim :: bool
where
  "hormander_euclidean_l2_fourier_plancherel_claim \<longleftrightarrow>
    (\<exists>fourier_l2 :: hormander_fourier_field \<Rightarrow>
        hormander_fourier_field.
      (\<forall>f. hormander_fourier_l2_on_plane f \<longrightarrow>
        hormander_fourier_l2_on_plane (fourier_l2 f)) \<and>
      (\<forall>f. hormander_fourier_l2_on_plane f \<and> integrable lborel f
        \<longrightarrow> (AE xi in lborel.
          fourier_l2 f xi = hormander_fourier_l1_transform f xi)) \<and>
      (\<forall>f. hormander_fourier_l2_on_plane f \<longrightarrow>
        integral\<^sup>L lborel (\<lambda>xi. norm (fourier_l2 f xi) ^ 2) =
          (2 * pi) ^ 2 *
            integral\<^sup>L lborel (\<lambda>x. norm (f x) ^ 2)) \<and>
      (\<forall>f. hormander_fourier_l2_on_plane f \<longrightarrow>
        (AE x in lborel.
          fourier_l2 (fourier_l2 f) x =
            of_real ((2 * pi) ^ 2) * f (-x))))"
\end{lstlisting}
% CF-ISABELLE-END: HORMANDER-EUCLIDEAN-L2-FOURIER-PLANCHEREL

\reviewlabel{Independent English translation of the Isabelle code}
% CF-BACK-TRANSLATION-BEGIN: HORMANDER-EUCLIDEAN-L2-FOURIER-PLANCHEREL
The claim chooses an $L^2$ Fourier representative on complex fields over $\mathbb R^2$. On its stated domain it has the exact unnormalised convention, Plancherel factor and inversion factor; no values are constrained outside that domain.
% CF-BACK-TRANSLATION-END: HORMANDER-EUCLIDEAN-L2-FOURIER-PLANCHEREL
% CF-LITERATURE-ENTRY-END: HORMANDER-EUCLIDEAN-L2-FOURIER-PLANCHEREL

% CF-LITERATURE-ENTRY-BEGIN: HORMANDER-QUADRATIC-STATIONARY-PHASE
\subsubsection{\texttt{\detokenize{HORMANDER-QUADRATIC-STATIONARY-PHASE}}}
\noindent\textit{Formal identifier:} \texttt{\detokenize{hormander_quadratic_stationary_phase_decay_claim}}.\par
\noindent\textit{Relationship to source:} derived.

\reviewlabel{Source attribution and locator}
% CF-SOURCE-ATTRIBUTION-BEGIN: HORMANDER-QUADRATIC-STATIONARY-PHASE/SOURCE-1
Lars Hörmander, The Analysis of Linear Partial Differential Operators I, first edition, Springer-Verlag, 1983, Lemma 7.7.3, equations (7.7.7)-(7.7.8), printed pp. 218-219
% CF-SOURCE-ATTRIBUTION-END: HORMANDER-QUADRATIC-STATIONARY-PHASE/SOURCE-1

\reviewlabel{Quoted source theorem statement}
% CF-SOURCE-STATEMENT-BEGIN: HORMANDER-QUADRATIC-STATIONARY-PHASE/SOURCE-1
\begin{sourcequote}
H\"ormander Lemma 7.7.3 assumes a symmetric nondegenerate matrix $A$ with $\operatorname{Im}A\geq0$. For every integer $k>0$, integer $s>n/2$, and Schwartz amplitude $u$, equation (7.7.7) bounds the oscillatory integral $\int u(x)e^{i\omega\langle Ax,x\rangle/2}\,dx$ minus its determinant factor times the finite expansion $T_k(\omega)$ by a constant times $|\omega|^{-n/2-k}$ and finitely many $L^2$ norms of derivatives of $u$; equation (7.7.8) defines that expansion.
\end{sourcequote}
% CF-SOURCE-STATEMENT-END: HORMANDER-QUADRATIC-STATIONARY-PHASE/SOURCE-1

\reviewlabel{Formalized derived mathematical result}
% CF-DERIVED-RESULT-BEGIN: HORMANDER-QUADRATIC-STATIONARY-PHASE
For a real symmetric matrix with trivial kernel and a compactly supported smooth integrable amplitude in any finite positive dimension, take $k=1$. The determinant main term decays like $\omega^{-n/2}$ and the remainder faster, so the displayed quadratic oscillatory integral tends to zero as the real parameter tends to positive infinity.
% CF-DERIVED-RESULT-END: HORMANDER-QUADRATIC-STATIONARY-PHASE

\reviewlabel{Isabelle code}
% CF-ISABELLE-BEGIN: HORMANDER-QUADRATIC-STATIONARY-PHASE
\begin{lstlisting}
definition hormander_compact_smooth_amplitude ::
  "(real^'n::finite \<Rightarrow> complex) \<Rightarrow> bool"
where
  "hormander_compact_smooth_amplitude u \<longleftrightarrow>
    smooth_on UNIV u \<and>
    compact (closure {x. u x \<noteq> 0}) \<and>
    integrable lborel u"

definition hormander_real_symmetric_matrix ::
  "(real^'n::finite^'n) \<Rightarrow> bool"
where
  "hormander_real_symmetric_matrix A \<longleftrightarrow>
    (\<forall>i j. A $ i $ j = A $ j $ i)"

definition hormander_real_nondegenerate_matrix ::
  "(real^'n::finite^'n) \<Rightarrow> bool"
where
  "hormander_real_nondegenerate_matrix A \<longleftrightarrow>
    (\<forall>x. A *v x = 0 \<longrightarrow> x = 0)"

definition hormander_quadratic_oscillatory_integral ::
  "(real^'n::finite^'n) \<Rightarrow> (real^'n \<Rightarrow> complex) \<Rightarrow>
    real \<Rightarrow> complex"
where
  "hormander_quadratic_oscillatory_integral A u omega =
    integral\<^sup>L lborel
      (\<lambda>x. exp (\<i> * of_real
        (omega * inner x (A *v x) / 2)) * u x)"

definition hormander_quadratic_stationary_phase_decay_claim ::
  "'n::finite itself \<Rightarrow> bool"
where
  "hormander_quadratic_stationary_phase_decay_claim dimension_type \<longleftrightarrow>
    (\<forall>(A :: real^'n^'n) u.
      hormander_real_symmetric_matrix A \<and>
      hormander_real_nondegenerate_matrix A \<and>
      hormander_compact_smooth_amplitude u
      \<longrightarrow>
      ((\<lambda>omega. hormander_quadratic_oscillatory_integral A u omega)
        \<longlongrightarrow> 0) at_top)"
\end{lstlisting}
% CF-ISABELLE-END: HORMANDER-QUADRATIC-STATIONARY-PHASE

\reviewlabel{Independent English translation of the Isabelle code}
% CF-BACK-TRANSLATION-BEGIN: HORMANDER-QUADRATIC-STATIONARY-PHASE
The formal decay claim universally quantifies the matrix and amplitude. Symmetry, trivial kernel, smooth compact support and integrability imply that the totalized Lebesgue integral with phase $e^{i\omega\langle Ax,x\rangle/2}$ has limit zero at the top filter.
% CF-BACK-TRANSLATION-END: HORMANDER-QUADRATIC-STATIONARY-PHASE
% CF-LITERATURE-ENTRY-END: HORMANDER-QUADRATIC-STATIONARY-PHASE
% CF-LITERATURE-ENTRIES-END

\clearpage
\section{Bibliography}
% CF-BIBLIOGRAPHY-BEGIN
% CF-BIBLIOGRAPHY-ENTRY: AIM-PLANAR-BEURLING-LP-BOUND
\subsection*{\texttt{\detokenize{AIM-PLANAR-BEURLING-LP-BOUND}}}
Kari Astala, Tadeusz Iwaniec, and Gaven Martin, Elliptic Partial Differential Equations and Quasiconformal Mappings in the Plane, Princeton Mathematical Series 48, Princeton University Press, 2009, formula (4.17), Theorem 4.0.10, and Theorem 4.5.3 with (4.84), printed pp. 94, 97-98, and 129, ISBN 978-0-691-13777-3
% CF-BIBLIOGRAPHY-ENTRY-END: AIM-PLANAR-BEURLING-LP-BOUND

% CF-BIBLIOGRAPHY-ENTRY: AIM-PLANAR-CAUCHY-BEURLING-DERIVATIVES
\subsection*{\texttt{\detokenize{AIM-PLANAR-CAUCHY-BEURLING-DERIVATIVES}}}
Kari Astala, Tadeusz Iwaniec, and Gaven Martin, Elliptic Partial Differential Equations and Quasiconformal Mappings in the Plane, Princeton Mathematical Series 48, Princeton University Press, 2009, Definition 4.0.9 and formulas (4.6), (4.8), (4.16)-(4.17), Theorem 4.0.10, the Lp extension discussion on printed pp. 96-97, and Theorem 4.5.3, printed pp. 93-98 and 129, ISBN 978-0-691-13777-3
% CF-BIBLIOGRAPHY-ENTRY-END: AIM-PLANAR-CAUCHY-BEURLING-DERIVATIVES

% CF-BIBLIOGRAPHY-ENTRY: AIM-PLANAR-CAUCHY-TEST-LEFT-INVERSE
\subsection*{\texttt{\detokenize{AIM-PLANAR-CAUCHY-TEST-LEFT-INVERSE}}}
Kari Astala, Tadeusz Iwaniec, and Gaven Martin, Elliptic Partial Differential Equations and Quasiconformal Mappings in the Plane, Princeton Mathematical Series 48, Princeton University Press, 2009, Definition 4.0.9 and formulas (4.6)-(4.8), printed p. 93 / PDF p. 112, ISBN 978-0-691-13777-3
% CF-BIBLIOGRAPHY-ENTRY-END: AIM-PLANAR-CAUCHY-TEST-LEFT-INVERSE

% CF-BIBLIOGRAPHY-ENTRY: AIM-PLANAR-HARDY-LITTLEWOOD-SOBOLEV
\subsection*{\texttt{\detokenize{AIM-PLANAR-HARDY-LITTLEWOOD-SOBOLEV}}}
Kari Astala, Tadeusz Iwaniec, and Gaven Martin, Elliptic Partial Differential Equations and Quasiconformal Mappings in the Plane, Princeton Mathematical Series 48, Princeton University Press, 2009, Theorem 4.3.8, printed p. 111; normalization and HLS estimate printed pp. 109-110
% CF-BIBLIOGRAPHY-ENTRY-END: AIM-PLANAR-HARDY-LITTLEWOOD-SOBOLEV

% CF-BIBLIOGRAPHY-ENTRY: AIM-PLANAR-RIESZ-POTENTIAL-HLS
\subsection*{\texttt{\detokenize{AIM-PLANAR-RIESZ-POTENTIAL-HLS}}}
Kari Astala, Tadeusz Iwaniec, and Gaven Martin, Elliptic Partial Differential Equations and Quasiconformal Mappings in the Plane, Princeton Mathematical Series 48, Princeton University Press, 2009, displayed Hardy-Littlewood-Sobolev inequality on printed p. 110 / PDF p. 129, ISBN 978-0-691-13777-3
% CF-BIBLIOGRAPHY-ENTRY-END: AIM-PLANAR-RIESZ-POTENTIAL-HLS

% CF-BIBLIOGRAPHY-ENTRY: EVANS-COMPACT-SMOOTH-L1-DENSITY
\subsection*{\texttt{\detokenize{EVANS-COMPACT-SMOOTH-L1-DENSITY}}}
Lawrence C. Evans, Partial Differential Equations, second edition, Graduate Studies in Mathematics 19, American Mathematical Society, 2010, Appendix C.5, standard mollifier printed p. 713 and Theorem 7(i),(iv), printed pp. 714-716
% CF-BIBLIOGRAPHY-ENTRY-END: EVANS-COMPACT-SMOOTH-L1-DENSITY

% CF-BIBLIOGRAPHY-ENTRY: EVANS-COMPACT-SMOOTH-LP-DENSITY
\subsection*{\texttt{\detokenize{EVANS-COMPACT-SMOOTH-LP-DENSITY}}}
Lawrence C. Evans, Partial Differential Equations, 2nd ed., Graduate Studies in Mathematics 19, American Mathematical Society, 2010, Appendix C.5 standard mollifier and Theorem 7(i),(iv), printed pp. 713-716, ISBN 978-0-8218-4974-3
% CF-BIBLIOGRAPHY-ENTRY-END: EVANS-COMPACT-SMOOTH-LP-DENSITY

% CF-BIBLIOGRAPHY-ENTRY: EVANS-COMPACT-SUPPORT-W1P-ZERO-DENSITY
\subsection*{\texttt{\detokenize{EVANS-COMPACT-SUPPORT-W1P-ZERO-DENSITY}}}
Lawrence C. Evans, Partial Differential Equations, second edition, Graduate Studies in Mathematics 19, American Mathematical Society, 2010, definition of C\_c\textasciicircum{}infinity(U), printed p. 256; definition of W\_0\textasciicircum{}\{k,p\}(U), printed p. 259; Section 5.3.1, Theorem 1 and proof, printed pp. 264-265; Appendix C.5, standard mollifier and Theorem 7(i),(iv), printed pp. 713-716; ISBN 978-0-8218-4974-3
% CF-BIBLIOGRAPHY-ENTRY-END: EVANS-COMPACT-SUPPORT-W1P-ZERO-DENSITY

% CF-BIBLIOGRAPHY-ENTRY: HORMANDER-EUCLIDEAN-L2-FOURIER-PLANCHEREL
\subsection*{\texttt{\detokenize{HORMANDER-EUCLIDEAN-L2-FOURIER-PLANCHEREL}}}
Lars Hormander, The Analysis of Linear Partial Differential Operators I: Distribution Theory and Fourier Analysis, first edition, Springer-Verlag, 1983, Definition 7.1.1, formulas (7.1.3) and (7.1.8), Definition 7.1.9, and Theorems 7.1.10-7.1.11, ISBN 3-540-12104-8 / 0-387-12104-8, DOI 10.1007/978-3-642-96750-4.
% CF-BIBLIOGRAPHY-ENTRY-END: HORMANDER-EUCLIDEAN-L2-FOURIER-PLANCHEREL

% CF-BIBLIOGRAPHY-ENTRY: HORMANDER-QUADRATIC-STATIONARY-PHASE
\subsection*{\texttt{\detokenize{HORMANDER-QUADRATIC-STATIONARY-PHASE}}}
Lars Hörmander, The Analysis of Linear Partial Differential Operators I, first edition, Springer-Verlag, 1983, Lemma 7.7.3, equations (7.7.7)-(7.7.8), printed pp. 218-219
% CF-BIBLIOGRAPHY-ENTRY-END: HORMANDER-QUADRATIC-STATIONARY-PHASE
% CF-BIBLIOGRAPHY-END
\end{document}
